
\begin{table}[H]
\centering
\caption{Structure-Activity Relationship Rules for SMVT Binding}
\label{tab:sar_rules}
\begin{tabular}{p{0.12\textwidth}p{0.35\textwidth}p{0.40\textwidth}}
\toprule
\textbf{Rule} & \textbf{Description} & \textbf{Structural Basis} \\
\midrule
R1: Barbiturate core & Malonylurea (O=C1CC(=O)NC(=O)N1) is the strongest pharmacophore. 8/8 barbiturates cross $-$7.0 kcal/mol & Mimics biotin's ureido ring; 2 N--H + 3 C=O in planar arrangement \\
\addlinespace
R2: Cyclic > linear & Cyclic ureide/amide required; linear analogs show weak binding & Pre-organized conformation reduces entropic penalty \\
\addlinespace
R3: Carboxyl dispensable & $-COOH$ not required (barbiturates lack it); but adds 0.5--1.0 kcal/mol when present & Salt bridge with Na$^+$ binding site \\
\addlinespace
R4: Size: 150--350 Da & Optimal MW range; >500 Da compounds fail (statins, bile acids) & Binding cavity accommodates $\sim$20 heavy atoms \\
\addlinespace
R5: logP: 0.5--3.5 & Balanced hydrophilicity; logP < 0 or > 4 reduces affinity & Membrane-adjacent pocket favors moderate lipophilicity \\
\addlinespace
R6: HBD: 2--4 & 2--4 hydrogen bond donors optimal; barbiturates have exactly 2 N--H & Orients ligand via H-bonds to backbone carbonyls \\
\addlinespace
R7: Planar pharmacophore & The N--C(=O)--N motif must be planar; non-planar analogs inactive & Required for simultaneous H-bond donor/acceptor presentation \\
\addlinespace
R8: Aromatic stacking & Phenyl at C5 enhances affinity (phenobarbital $-8.30$ > butalbital $-7.73$) & $\pi$-stacking with aromatic residues in hydrophobic pocket \\
\addlinespace

\bottomrule
\end{tabular}
\end{table}
